\subsection{Two candidate first primes remove the quintic spin assumption}

\begin{result}{E.213}{Proved}
{A legal original first pair with invertible five-primary first matrix exists}
There are legal first primes $q_0,r_0$ for the unchanged original
factorial recursion E.102, with the minimal five-adic shells of
E.199, for which the matrix in E.202.1 is invertible.
Here \emph{legal} includes the entire original finite set of old
local congruences, the original radical-splitting conditions,
the current-prime congruence at $q_0$, and every original lower
bound and distinctness requirement. No bound is placed on how
large the resulting primes may be.

Consequently the original family has an ordinary infinite-support
member whose unique integral pro-five class-two section matches
the ordinary local point class at every support prime. Its
continuous pro-five class-three obstruction vanishes, and it has
an integral pro-five class-three lift. This consequence uses
E.211; it supplies neither class-three support matching nor a
full pro-five or full pronilpotent section.

The proof does not prescribe the quintic spin of a varying
prime. It first chooses two legal candidate $A$-primes and
then one prime which is a legal $B$-prime for both candidates.
Their common second diagonal cannot make both matrices singular.
\end{result}

\begin{proof}
\emph{Fix the complete original first-stage local data.}
At stage zero the exponent is $24$, the next exponent is $120$,
and the initial coefficients are $3,-2$. Fix once and for all
the exact old unit roots used by E.102, with E.199's minimal
five-adic shell. Thus all old $A$-prime congruences are fixed,
all old $B$-prime congruences are fixed, and the old set contains
every prime up to
\[
                         B(120)=121^4=214358881.
 \tag{E.213.1}
\]
The precision is the actual precision of E.102, not merely a
congruence modulo $120$ or $25$.

The first-prime Chebotarev step of E.102 supplies infinitely
many eligible primes for these fixed old $A$-coordinates.
Choose two distinct ones, denoted $q^{(1)},q^{(2)}$, both
larger than every original exclusion. Each satisfies all the
original first-prime conditions, including complete splitting
in
\[
       \mathbb Q\bigl(\mu_{24},(-2)^{1/24}\bigr)
 \quad\text{and}\quad q^{(i)}\equiv1\pmod{120}.
 \tag{E.213.2}
\]
Their prescribed five-adic shells have
$v_5(q^{(i)}-1)=1$. The two candidates share the same old
local data; no relation between their unknown self residues
has been assumed.

For the fixed local frames of E.202 set
\[
 a_i=\ell_{q^{(i)}}\bigl(3U_{q^{(i)}}\bigr)\in\mathbb F_5.
 \tag{E.213.3}
\]
If either $a_i$ is zero, E.208 already gives a legal second
prime with invertible first matrix. We may therefore assume
throughout the remaining argument that $a_1,a_2\ne0$.

\emph{Four independently variable fixed-field coordinates.}
Let $K=\mathbb Q(\zeta_5)$. For $i=1,2$ use the classes
$z_i=z_{q^{(i)}}$ of E.202 and $w_i=w_{q^{(i)}}$ of E.206.
They are defined before the second prime is chosen. Each has
tame residue one at its own $q^{(i)}$ and zero at the other
candidate; $w_i$ also has zero localization at five.
Their restrictions to $G_K$ cut out cyclic degree-five
extensions, Galois over $\mathbb Q$, with respective
conjugation characters $\chi^2,\chi^3$.

Within either character the two extensions are independent:
a relation would have zero residue at $q^{(1)}$ and
$q^{(2)}$, forcing both coefficients to vanish.
The two different characters are independent because the
order-four group $\operatorname{Gal}(K/\mathbb Q)$ acts
semisimply over $\mathbb F_5$. Their joint field $L$ therefore
satisfies
\[
 \operatorname{Gal}(L/K)
   =\mathbb F_5(2)^2\oplus\mathbb F_5(3)^2.                 \tag{E.213.4}
\]
Its maximal abelian subextension over $\mathbb Q$ is exactly
$K$: the commutator contains every nontrivial-character
factor in E.213.4. In particular, for every cyclotomic
field $F$ containing $K$,
\[
 L\cap F=K,\qquad
 \operatorname{Gal}(LF/F)
   =\mathbb F_5(2)^2\oplus\mathbb F_5(3)^2.                 \tag{E.213.5}
\]
These are independent values of four fixed cohomology classes,
not four conditions on the self residue of the future prime.

\emph{One second prime can retain both original radical conditions.}
Put
\[
 K_0=\mathbb Q(\mu_{120}),\qquad
 A_i=3(q^{(i)})^{24},\qquad
 \theta_i=A_i^{1/120}>0,\qquad
 R_0=K_0(\theta_1,\theta_2).
 \tag{E.213.6}
\]
All roots are chosen compatibly in an algebraic closure.
Let $u=3^{1/24}>0$. The exact identities
\[
 \theta_i^5=q^{(i)}u,\qquad
 \theta_i^{24}=A_i^{1/5},\qquad
 \theta_i=\frac{(q^{(i)}u)^5}{A_i^{1/5}}
 \tag{E.213.7}
\]
give
\[
 R_0=K_0\bigl(u,A_1^{1/5},A_2^{1/5}\bigr).
 \tag{E.213.8}
\]
Thus the non-five part of the two compulsory radical fields
is the same field, whereas their five-primary parts are
ordinary weight-one Kummer lines.

For completeness, $[K_0(u):K_0]=12$. The square root of
three lies in $K_0$, so this degree divides $12$.
If its two-primary part were smaller than four, then
$u^6=3^{1/4}$ would lie in the abelian extension $K_0$.
If its three-primary part were trivial, then
$u^8=3^{1/3}$ would lie there. Both are impossible because
the real fields generated by these irreducible pure roots
are not normal over $\mathbb Q$, whereas every subfield of
an abelian Galois extension is normal.
The two classes $A_1,A_2$ are independent modulo fifth powers
over $K_0$: at a prime above $q^{(i)}$ the valuation of $A_i$
is $24$, prime to five, and that of the other class is zero.
These primes are unramified in $K_0$. The degree-twelve
extension does not alter this independence. Consequently
\[
 N:=\operatorname{Gal}(R_0/K_0)
       \simeq C_{12}\times C_5^2.                            \tag{E.213.9}
\]

We need the maximal abelian subextension of $R_0/\mathbb Q$,
not just the degrees of its factors. The field $R_0$ is
Galois over $\mathbb Q$ and $K_0/\mathbb Q$ is abelian.
Complex conjugation acts on every Kummer translation in
$N$ by inversion, so
\[
       2N\subseteq
       [\operatorname{Gal}(R_0/\mathbb Q),
        \operatorname{Gal}(R_0/\mathbb Q)]\subseteq N.
 \tag{E.213.10}
\]
The quotient $N/2N$ has order two. Its fixed field is
\[
                         R_0^{2N}=K_0(3^{1/4}).              \tag{E.213.11}
\]
This field is not abelian over $\mathbb Q$, since it
contains the nonnormal subfield $\mathbb Q(3^{1/4})$.
The commutator in E.213.10 therefore cannot equal $2N$;
it equals $N$. Hence the maximal abelian subextension of
$R_0/\mathbb Q$ is precisely $K_0$.

At each candidate $q^{(i)}$ the original splitting condition
in E.213.2 and simple-root lifting give a local unit $t_i$
with
\[
       -2t_i^{24}=1-3(q^{(i)})^{24}
       \quad\text{in }\mathbb Q_{q^{(i)}}.                   \tag{E.213.12}
\]
Add the two simultaneous original current-prime congruences
$r\equiv t_i\pmod{q^{(i)}}$ to the common old $B$-coordinate
congruences. Their moduli are relatively prime to each
other and to the old moduli. The resulting CRT class is a
unit class and is still one modulo $120$.
Let $F$ be the full cyclotomic CRT field, containing
$K_0$, and let $\sigma_F$ be this automorphism of $F$.
It is the identity on $K_0$.

Since $R_0\cap F=K_0$, the maps $\sigma_F$ on $F$ and
the identity on $R_0$ glue to an automorphism of $FR_0$.
It fixes both chosen $120$th roots. This proves actual
simultaneous compatibility of the two original radical
requirements with every old congruence and both new
current-prime congruences. In particular
\[
 \operatorname{Gal}(FR_0/F)\simeq C_{12}\times C_5^2,
 \tag{E.213.13}
\]
whose exponent-five quotient has conjugation character
$\chi$ on both factors.

\emph{The four extra coordinates do not disturb that compatibility.}
The intersection $LF\cap FR_0$ is Galois over $\mathbb Q$.
Its group over $F$ is an elementary abelian five-group,
since it is a quotient of E.213.5. By E.213.13 it is also
a quotient of two copies of the character $\chi$.
But E.213.5 has only characters $\chi^2,\chi^3$.
The three characters are distinct over $\mathbb F_5$;
therefore this common quotient is zero:
\[
                             LF\cap FR_0=F.                  \tag{E.213.14}
\]
The argument remains valid for the entire CRT field,
including its cyclotomic five-power part: the action on
all the displayed five-groups factors through the same
order-four cyclotomic quotient.

Because $L\cap F=K$ and $\sigma_F$ is trivial on $K$,
we may prescribe its lift on $LF$ with arbitrary values
of $z_1,w_1,z_2,w_2$. By E.213.14 this lift glues to
the just-constructed automorphism of $FR_0$.
Choose
\[
 z_1(F_r)=z_2(F_r)=1,\qquad
 w_1(F_r)=a_1,\qquad w_2(F_r)=2a_2.                           \tag{E.213.15}
\]
At this point these symbols specify an automorphism of
the finite joint field; the prime $r$ has not yet been
chosen.

Chebotarev supplies infinitely many rational primes in
its conjugacy class. Conjugation multiplies a $z_i$ value
by $\chi^2$ and a $w_i$ value by $\chi^3$, since the
prescribed Frobenius has $\chi=1$ modulo five.
Thus each product $z_i(F_r)w_i(F_r)$ is multiplied by
the same scalar $\chi^5=\chi$. In particular the two
ratios
\[
      \frac{z_1(F_r)w_1(F_r)}{a_1},
      \qquad
      \frac{z_2(F_r)w_2(F_r)}{a_2}                            \tag{E.213.16}
\]
remain distinct under conjugation. All cyclotomic CRT
conditions are invariant under conjugacy, and restriction
to $R_0$ remains the identity. Choose such a prime $r$
larger than both candidates and all original exclusions.
It is a legal original second prime for either candidate,
with the original minimal five-adic shell. No selection
using the value of its own diagonal has occurred.

\emph{The same second diagonal cannot make both matrices singular.}
Use one local frame at $r$, common to both prospective
first pairs. The second diagonal in both E.202 matrices
is exactly the same element
\[
                         d=\ell_r(-2U_r)\in\mathbb F_5.
 \tag{E.213.17}
\]
The reciprocity identities E.206 identify their other
entries as
\[
 M_i=
 \begin{pmatrix}
 a_i&w_i(F_r)\\ z_i(F_r)&d
 \end{pmatrix},
 \qquad
 \det M_i=a_i d-z_i(F_r)w_i(F_r).
 \tag{E.213.18}
\]
Since $a_i\ne0$, the two determinants could vanish
simultaneously only if $d$ equalled both distinct ratios
in E.213.16. This is impossible. Select a candidate for
which the determinant is nonzero and call it $q_0$;
put $r_0=r$. The unused candidate merely contributed
additional finite congruences to the prime selection.
All requirements of the original recursion for the
selected pair still hold.

E.202 now has a unique successful parameter in its
twenty-five allowed stage-one choices. E.211 continues
that parameter inside the unchanged original recursion,
matching the integral pro-five class-two section at all
support primes. Its class-three existence consequence
proves the assertion. No quintic-spin distribution
statement or varying-prime Chebotarev condition has been
used.

The finite certificate
\path{code/ghost_audit_two_first_primes.py} checks the
two-matrix argument for every pair of nonzero first
diagonals, every second diagonal and every cyclotomic
conjugating scalar, as well as the distinct-character
quotient calculation and the elementary radical-group
identities. Its output is
\path{code/ghost_audit_two_first_primes.out}.
It is not a numerical production of primes satisfying
the original large CRT conditions; their existence is
the finite-field-intersection and Chebotarev proof above.
\end{proof}

\paragraph{Refinement: every later weight-three localization can be a unit.}
The construction can additionally require
\[
 \ell_{q_j}(\epsilon)\ne0,\qquad
 \ell_{r_j}(\epsilon)\ne0
 \quad\text{for every }j\ge0,\qquad
 \epsilon=(1+\sqrt5)/2,
 \tag{E.213.19}
\]
with the Kummer and Tate frames of E.206.
This is a simultaneous refinement of the proved construction,
not an assumption on the unknown spin or on the next
class-three discrepancies.

Indeed, let $T/K$ be the degree-five extension cut out by
the global weight-three epsilon class. E.206 proves that
this class is unramified outside five and is nonzero
locally at five. The field is Galois over $\mathbb Q$,
with conjugation character $\chi^3$, and its maximal
abelian subextension over $\mathbb Q$ is $K$.
At the first $A$-prime step the compulsory radical field
$\mathbb Q(\mu_{24},(-2)^{1/24})$ has degree prime to
five. After adjoining the entire cyclotomic CRT field
$F$, it is therefore disjoint from $TF$ over $F$.
Since the prescribed CRT automorphism is the identity
on $K$, either candidate $q^{(i)}$ may be chosen with
arbitrary nonzero epsilon coordinate, retaining every
original congruence and radical splitting.

At the common $B$-prime step, $T$ is independent of
the four fields in E.213.4. The weight-two constituents
have a different character. A relation among the three
weight-three constituents is zero at five only if its
epsilon coefficient is zero, since $w_1,w_2$ were
normalized to be zero there. The two remaining
coefficients vanish by their private residues.
Consequently adjoining $T$ replaces E.213.5 by
\[
 \operatorname{Gal}(LTF/F)
   =\mathbb F_5(2)^2\oplus\mathbb F_5(3)^3.
 \tag{E.213.20}
\]
It still has no weight-one quotient, so the exact
intersection argument E.213.14 is unchanged. We may
prescribe a nonzero epsilon coordinate as well as
E.213.15. Conjugation preserves its nonvanishing and
the inequality of the two ratios. If an initial
diagonal is zero, the identical argument with just
that candidate augments E.208. Thus the legal
selected pair has both epsilon values nonzero.

At any subsequent prime selection, E.205 combines
the original compulsory radical over the full
cyclotomic CRT field $F$ with one independent
degree-five rational Kummer line: the $A_0$ line
for a new $q_j$, or the $B_0$ line for a new $r_j$.
Those existing field-intersection arguments retain
the prescribed fixed root and give the required
nonzero $\alpha$ or $\beta$. This entire compositum
is abelian over $F$, and every exponent-five
quotient of its Galois group has character $\chi$:
all its generators are roots of rational numbers.
It is therefore disjoint from $TF/F$, whose
character is $\chi^3$. The epsilon Frobenius value
can be prescribed nonzero independently.
This conclusion holds for every full CRT field
and every one of its prescribed automorphisms
in the recursion. In particular it does not
restrict any of the old root-of-unity choices
used by the exact affine solve E.211. All
principal-root retention, actual ordinary
limits, and class-two support matches proved
there remain valid.

Normalize the integral global weight-three
generator $\zeta_3$ to reduce to the epsilon
class, as in E.212. Then E.213.19 makes each
$\tau_p$ at the fresh support a unit. For any
one chosen global class-three lift, every
individual local matching preimage is nonempty
and is exactly
\[
 C_p=\tau_p^{-1}d_p+
            5^{e_j}\mathbb Z_5^2
 \quad(p=q_j\text{ or }r_j).
 \tag{E.213.21}
\]
The support places two and three have zero
local weight-three correction group.
The remaining class-three support gate is
therefore precisely the pairwise compatibility
of the two-vector residues in E.213.21.
Since $e_j\longrightarrow\infty$, at most one
global correction can meet all these tests.
The refinement does not establish their
compatibility, and it imposes no point match
at the coefficient place five.
